\section{Methodology}
\label{sec:method}

This section describes the evaluation protocol shared by all experiments and then the three modelling approaches.
\Cref{fig:architecture} gives an overview of how the data flows through the project.

\begin{figure}[H]
  \centering
  \begin{tikzpicture}[
      font=\footnotesize,
      box/.style={draw, rounded corners, align=center, minimum height=1.1cm, fill=#1},
      box/.default=gray!12,
      arr/.style={-{Stealth[length=2mm]}, thick}
    ]
    % Row 1: shared data preparation
    \node[box, text width=3.4cm] (raw) at (-5, 0) {Raw CSV\\Kaggle, 32{,}581 loans};
    \node[box, text width=3.4cm] (clean) at (0, 0) {Rule-based cleaning\\duplicates, impossible values};
    \node[box, text width=3.4cm] (split) at (5, 0) {Stratified split\\80\% train / 20\% test};
    \draw[arr] (raw) -- (clean);
    \draw[arr] (clean) -- (split);
    % Row 2: the three model families
    \node[box=blue!8, text width=4.3cm, minimum height=2.2cm] (trad) at (-5, -3)
      {\textbf{Traditional pipeline}\\impute, engineer, clip, encode, scale, select / PCA, resample, classifier (8 models)};
    \node[box=orange!12, text width=4.3cm, minimum height=2.2cm] (dl) at (0, -3)
      {\textbf{Deep learning pipeline}\\impute and scale numerics, integer codes, entity embeddings, MLP / TabNet};
    \node[box=green!10, text width=4.3cm, minimum height=2.2cm] (bn) at (5, -3)
      {\textbf{Bayesian network}\\quantile bins, structure learning, CPTs, exact and approximate inference};
    \coordinate (bus) at (0, -1.25);
    \draw[thick] (split.south) |- (bus -| trad.north);
    \draw[thick] (bus -| split.south) -- (bus -| trad.north);
    \draw[arr] (bus -| trad.north) -- (trad.north);
    \draw[arr] (bus -| dl.north) -- (dl.north);
    \draw[arr] (bus -| bn.north) -- (bn.north);
    % Row 3: shared evaluation and outputs
    \node[box, minimum width=14.2cm] (eval) at (0, -6.2)
      {\textbf{Common evaluation}: 5-fold stratified CV on the training set, decision threshold tuned for F2,\\
       one final test-set evaluation (PR-AUC, ROC-AUC, F2, recall, precision); features saved as \texttt{.npy} / \texttt{.h5}};
    \draw[arr] (trad.south) -- (trad.south |- eval.north);
    \draw[arr] (dl.south) -- (dl.south |- eval.north);
    \draw[arr] (bn.south) -- (bn.south |- eval.north);
  \end{tikzpicture}
  \caption{Overview of the project: one cleaned and split dataset feeds three model families that share the same evaluation protocol.}
  \label{fig:architecture}
\end{figure}

\subsection{Evaluation protocol}
\label{subsec:protocol}

\paragraph{Data split.}
After rule-based cleaning (\Cref{subsec:quality}), the data are split once into a stratified training set (80\%) and test set (20\%) with a fixed seed.
Every learned step (imputation statistics, scaling parameters, PCA axes, discretization bins, network weights, CPTs, decision thresholds) is estimated on training data only.
The test set is used once per final model.

\paragraph{Cross-validation.}
Model selection uses 5-fold stratified cross-validation on the training set.
Because every preprocessing step lives inside one pipeline object, it is refitted on each training fold, so the validation fold never leaks into the preprocessing statistics.

\paragraph{Metrics.}
Let $TP, FP, FN, TN$ count the outcomes for the positive class (default).
We report
\begin{align}
  \text{precision} &= \frac{TP}{TP + FP}, &
  \text{recall} &= \frac{TP}{TP + FN}, &
  \Fbeta &= \frac{(1 + \beta^2)\,\text{precision}\cdot\text{recall}}{\beta^2\,\text{precision} + \text{recall}},
  \label{eq:fbeta}
\end{align}
together with accuracy, the area under the ROC curve (ROC-AUC) and the area under the precision-recall curve (PR-AUC, computed as average precision).

\paragraph{Why PR-AUC and F2 instead of accuracy or recall alone.}
Defaults make up only 21.8\% of the loans, so the trivial ``never default'' rule already reaches 78\% accuracy.
For a lender, a missed defaulter (false negative) loses the principal, whereas a rejected good applicant (false positive) loses only the interest margin.
Recall on the default class is therefore the quantity of interest, but recall alone is degenerate: flagging every applicant gives recall 1.
We consequently
\begin{enumerate}[nosep]
  \item rank models and select hyper-parameters by \textbf{PR-AUC}, a threshold-free measure that, unlike ROC-AUC, is sensitive to performance on the minority class \cite{saito2015}; and
  \item choose each model's operating threshold to maximize \textbf{F2} ($\beta = 2$ in \cref{eq:fbeta}), which weights recall twice as much as precision \cite{rijsbergen1979}.
\end{enumerate}
The threshold is tuned by internal cross-validation on the training set (\ident{TunedThresholdClassifierCV}) for the traditional models, and on a held-out 15\% validation split of the training set for the neural networks and Bayesian networks.

\subsection{Traditional machine learning pipeline}
\label{subsec:trad}

The pipeline is assembled by a single function from a declarative \ident{PipelineConfig}, shown in \cref{lst:config}.
All components are standard scikit-learn \cite{sklearn} or imbalanced-learn \cite{imblearn} classes, except two small transformers that have no library equivalent: grade-wise median imputation and IQR clipping.

\begin{listing}[H]
  \begin{minted}[fontsize=\footnotesize, frame=lines, bgcolor=gray!5]{python}
@dataclass(frozen=True)
class PipelineConfig:
    imputer: str = "median"          # mean | median | group_median | knn | iterative
    add_indicator: bool = False      # append missing-value indicator columns
    feature_engineering: bool = False
    outlier: str = "none"            # none | clip_iqr | log
    grade_encoding: str = "ordinal"  # ordinal | onehot
    scaler: str = "standard"         # none | standard | minmax | robust
    feature_range: tuple = (0.0, 1.0)  # used by MinMaxScaler
    select_k: int | None = None      # SelectKBest(mutual information)
    pca_variance: float | None = None  # e.g. 0.90, 0.95
    resampler: str = "none"          # none | smote | undersample
    class_weight: str | None = None  # None | "balanced"
    model: str = "logreg"            # one of the 8 classifiers below
    model_params: dict = field(default_factory=dict)

pipe = build_pipeline(PipelineConfig(imputer="group_median", scaler="minmax",
                                     feature_range=(-1, 1), pca_variance=0.95, model="svm"))
  \end{minted}
  \caption{Every pipeline step is selected through one configuration object.}
  \label{lst:config}
\end{listing}

\subsubsection{Missing-value imputation}
Five strategies are available.
\emph{Mean} and \emph{median} imputation replace a missing value by a training statistic of its column.
\emph{Grade-wise median} imputation fills \ident{loan_int_rate} with the median rate of the applicant's loan grade, which the EDA shows to be the main driver of the rate (\cref{fig:rate_grade}).
\emph{KNN} imputation \cite{troyanskaya2001} averages the five nearest neighbours in standardized feature space.
\emph{Iterative} imputation (MICE) \cite{mice} regresses each incomplete column on the others in round-robin fashion.
Optionally, binary missing-value indicators are appended so that the model can exploit informative missingness \cite{rubin1987}.

\subsubsection{Feature engineering and outliers}
Four domain features can be added: $\log(1+\text{income})$, $\log(1+\text{amount})$, income per year of credit history, and the product of the interest rate and the loan-to-income ratio, which approximates the repayment burden.
Outliers are handled by IQR winsorizing, which clips every value $x$ of a column to $[Q_1 - 1.5\,\mathrm{IQR},\; Q_3 + 1.5\,\mathrm{IQR}]$ with quartiles learned on the training fold, or by a $\log(1+x)$ transform of all numeric columns.

\subsubsection{Encoding and scaling}
Nominal variables (\ident{person_home_ownership}, \ident{loan_intent}, \ident{cb_person_default_on_file}) are one-hot encoded.
The loan grade has a natural order A~$<$~B~$<\dots<$~G, so it can be encoded either ordinally (0--6) or one-hot.
Numeric features can be standardized ($z = (x - \mu)/\sigma$), min-max scaled to a configurable range $[a, b]$ via $x' = a + (b - a)(x - x_{\min})/(x_{\max} - x_{\min})$, or robust-scaled with the median and IQR.

\subsubsection{Feature selection and dimensionality reduction}
\ident{SelectKBest} keeps the $k$ encoded features with the highest mutual information with the target.
PCA \cite{jolliffe2016} projects the scaled features onto the leading principal components that together retain a chosen fraction of the variance (80\%, 90\%, 95\% or 99\%).

\subsubsection{Class imbalance}
Three remedies are compared: class weights inversely proportional to class frequency ($w_c = n / (2 n_c)$), SMOTE \cite{smote}, which synthesizes minority samples by interpolating between minority neighbours, and random under-sampling of the majority class.
Resampling is applied only to the training folds (it is a pipeline step that is skipped at prediction time).

\subsubsection{Classifiers}
\Cref{tab:classifiers} lists the eight classifiers.
Six of them are suggested by the assignment; Gaussian Naive Bayes and histogram-based gradient boosting are added as a probabilistic baseline and as a strong modern tree ensemble.

\begin{table}[H]
  \centering
  \small
  \begin{tabularx}{\linewidth}{lX}
    \toprule
    \textbf{Classifier} & \textbf{Decision function / idea} \\
    \midrule
    Logistic Regression & $P(y=1 \mid \bm{x}) = \sigma(\bm{w}^\top \bm{x} + b)$, $L_2$-regularized log-loss with strength $1/C$ \\
    Linear SVM & Hinge loss with a linear decision boundary \\
    SVM (RBF) & Maximum-margin classifier with kernel $k(\bm{x}, \bm{x}') = \exp(-\gamma \lVert \bm{x} - \bm{x}' \rVert^2)$ \cite{cortes1995} \\
    $k$-NN & Majority (or distance-weighted) vote of the $k$ nearest training points \\
    Gaussian Naive Bayes & $P(y \mid \bm{x}) \propto P(y) \prod_j \mathcal{N}(x_j; \mu_{jy}, \sigma^2_{jy})$ \\
    Decision Tree & Recursive axis-aligned splits minimizing Gini impurity \\
    Random Forest & Bagged decorrelated trees with random feature subsets \cite{breiman2001} \\
    Hist.\ Gradient Boosting & Additive trees fitted to the gradient of the log-loss on binned features \cite{friedman2001gbm,lightgbm} \\
    \bottomrule
  \end{tabularx}
  \caption{Classifiers available in the traditional pipeline.}
  \label{tab:classifiers}
\end{table}

\subsection{Deep learning pipeline}
\label{subsec:dl}

The deep pipeline replaces hand-made encodings with representations learned end to end (PyTorch \cite{pytorch}).
Preprocessing is reduced to what the optimizer needs: numeric columns are imputed, clipped and standardized, and each categorical column is mapped to integer codes, with code 0 reserved for categories unseen during training.

\paragraph{Entity-embedding MLP.}
Following \cite{entity_embeddings}, every categorical variable $c$ with $n_c$ levels owns a trainable embedding table $E_c \in \mathbb{R}^{n_c \times d_c}$ with $d_c = \min(8, \lceil n_c / 2 \rceil)$.
The embeddings are concatenated with the numeric vector and passed through $L$ blocks of Linear $\to$ BatchNorm \cite{batchnorm} $\to$ ReLU $\to$ Dropout \cite{dropout}, followed by a linear output producing a logit $z$.
Training minimizes the class-weighted binary cross-entropy
\begin{equation}
  \mathcal{L} = -\frac{1}{n}\sum_{i=1}^n \Big[ w_+ \, y_i \log \sigma(z_i) + (1 - y_i) \log\big(1 - \sigma(z_i)\big) \Big],
  \qquad w_+ = \frac{n_0}{n_1},
\end{equation}
with AdamW \cite{adamw} (learning rate $2\times10^{-3}$, weight decay $10^{-4}$, batch size 512), a learning rate halved when the validation PR-AUC stalls for three epochs, and early stopping with patience 10 on the validation PR-AUC.
Seven variants (\cref{tab:mlp}) vary the width, depth, dropout, embedding versus one-hot input, class weighting and the engineered features.
The 64-dimensional output of the last hidden layer is saved as a learned feature set.

\paragraph{TabNet.}
TabNet \cite{tabnet} processes the features in several decision steps.
At each step a sparse attention mask (sparsemax) selects which features to use, which yields built-in, instance-wise feature importance.
We use $n_d = n_a = 16$, four steps, relaxation $\gamma = 1.5$, 4-dimensional categorical embeddings and inverse-frequency sampling against the class imbalance.

\subsection{Bayesian network (extension for the talented-engineering program)}
\label{subsec:bn}

A Bayesian network (BN) \cite{koller2009} describes the joint distribution of discrete variables $X_1, \dots, X_n$ with two ingredients.
A directed acyclic graph $G$ encodes the conditional independences, and each node carries a conditional probability table (CPT) given its parents:
\begin{equation}
  P(X_1, \dots, X_n) = \prod_{i=1}^{n} P\big(X_i \mid \Pa_G(X_i)\big).
  \label{eq:bn_factorization}
\end{equation}
Everything below is implemented from scratch with NumPy in \ident{modules/creditrisk/bayesnet/}.
\Cref{tab:bn_modules} maps the components to files.

\begin{table}[H]
  \centering
  \small
  \begin{tabularx}{\linewidth}{lX}
    \toprule
    \textbf{File} & \textbf{Content} \\
    \midrule
    \ident{graph.py} & DAG with cycle detection on insertion, Kahn's topological sort, Markov blanket \\
    \ident{factor.py} & Discrete factor with product, marginalization (sum-out) and evidence reduction \\
    \ident{network.py} & CPT estimation, log-likelihood, vectorized ancestral sampling \\
    \ident{structure.py} & Naive Bayes, TAN (Chow-Liu), BIC hill climbing, expert DAG \\
    \ident{inference.py} & Variable elimination, enumeration, rejection sampling, likelihood weighting, Gibbs sampling \\
    \ident{discretize.py} & Quantile binning of numeric features learned on the training split \\
    \bottomrule
  \end{tabularx}
  \caption{Bayesian-network modules written for this project.}
  \label{tab:bn_modules}
\end{table}

\subsubsection{Discretization and missing values}
Each numeric feature is cut into $B$ quantile bins learned on the training split ($B$ is chosen on a validation split, \cref{tab:bn_sens}).
Missing values are \emph{not} imputed: they stay unobserved.
During learning, the CPT of $X_i$ is estimated from the rows in which $X_i$ and all its parents are observed (available-case analysis).
During prediction, a missing feature is simply left out of the evidence and marginalized by inference, which is the principled treatment of data missing at random.

\subsubsection{Parameter learning}
With a symmetric Dirichlet($\alpha$) prior, the maximum a posteriori CPT entries are smoothed relative frequencies
\begin{equation}
  \hat{\theta}_{x \mid \bm{u}} = \frac{N(x, \bm{u}) + \alpha}{N(\bm{u}) + \alpha\, r_i},
\end{equation}
where $N(\cdot)$ are counts in the training data and $r_i$ is the number of states of $X_i$.
$\alpha = 1$ is Laplace smoothing; the sensitivity to $\alpha$ is studied in \cref{tab:bn_sens}.

\subsubsection{Structures}
Four structures are compared.
\begin{itemize}[nosep]
  \item \textbf{Naive Bayes}: the class is the only parent of every feature.
  \item \textbf{TAN} (tree-augmented Naive Bayes) \cite{friedman1997bnc}: Naive Bayes plus a maximum-weight spanning tree over the features, where the weight of the pair $(X_i, X_j)$ is the conditional mutual information $I(X_i; X_j \mid C)$ \cite{chowliu1968}. The tree is built with Prim's algorithm.
  \item \textbf{Hill climbing on BIC}: greedy search over single-edge additions, deletions and reversals that keep the graph acyclic, with at most three parents per node. It maximizes
  \begin{equation}
    \mathrm{BIC}(G) = \sum_{i} \Big[ \sum_{x, \bm{u}} N(x, \bm{u}) \log \frac{N(x, \bm{u})}{N(\bm{u})} - \frac{\log N_i}{2} \, q_i (r_i - 1) \Big],
  \end{equation}
  where $q_i$ is the number of parent configurations \cite{schwarz1978}. Because the score decomposes over families, each move is re-scored locally and family scores are cached.
  \item \textbf{Expert DAG}: a causal story of a loan application written before inspecting the learned graphs. Age drives employment length, credit history and income. Income and amount determine the loan-to-income ratio. Past defaults determine the grade, which sets the interest rate. Default depends on the grade, the repayment burden, housing and the loan purpose.
\end{itemize}

\subsubsection{Exact inference: variable elimination}
Any query $P(Q \mid \bm{e})$ is answered by summing hidden variables out of the product of reduced CPT factors (\cref{alg:ve}) \cite{zhangpoole1994}.
Two optimizations matter.
First, variables that are not ancestors of $Q$ or $\bm{e}$ are \emph{barren}: their factors sum to one and are dropped before elimination.
Second, the elimination order is chosen greedily by the min-weight heuristic, which eliminates next the variable whose intermediate factor would be smallest.
For correctness, the unit tests compare variable elimination with brute-force enumeration of the full joint on random networks and require agreement to $10^{-12}$.

\begin{algorithm}[H]
  \caption{Variable elimination for $P(Q \mid \bm{e})$}
  \label{alg:ve}
  \KwIn{network $(G, \{\theta_i\})$, query variable $Q$, evidence $\bm{e}$}
  $R \leftarrow$ ancestors of $\{Q\} \cup \mathrm{vars}(\bm{e})$ in $G$ \tcp*{prune barren nodes}
  $\mathcal{F} \leftarrow \{\, \phi_i = P(X_i \mid \Pa(X_i))[\bm{e}] : X_i \in R \,\}$ \tcp*{reduce by evidence}
  $H \leftarrow R \setminus (\{Q\} \cup \mathrm{vars}(\bm{e}))$\;
  \For{$Z \in \textsc{MinWeightOrder}(\mathcal{F}, H)$}{
    $\mathcal{F}_Z \leftarrow \{\phi \in \mathcal{F} : Z \in \mathrm{scope}(\phi)\}$\;
    $\tau \leftarrow \sum_{z} \prod_{\phi \in \mathcal{F}_Z} \phi$\;
    $\mathcal{F} \leftarrow (\mathcal{F} \setminus \mathcal{F}_Z) \cup \{\tau\}$\;
  }
  \Return normalize$\big(\prod_{\phi \in \mathcal{F}} \phi\big)$\;
\end{algorithm}

\subsubsection{Approximate inference}
Three sampling algorithms are implemented and vectorized over samples \cite{aima,koller2009}.
\begin{itemize}[nosep]
  \item \textbf{Rejection sampling} draws $N$ samples in topological order from \cref{eq:bn_factorization} and keeps those consistent with $\bm{e}$. The acceptance rate equals $P(\bm{e})$, which shrinks exponentially with the number of evidence variables.
  \item \textbf{Likelihood weighting} \cite{fungchang1990} clamps the evidence variables, samples the others and weights each sample by $w = \prod_{E_j \in \bm{e}} P(e_j \mid \Pa(E_j))$; the estimate is $\hat{P}(q \mid \bm{e}) = \sum_s w_s \mathbb{1}[q_s = q] / \sum_s w_s$ (\cref{alg:lw}).
  \item \textbf{Gibbs sampling} \cite{geman1984} runs 20 parallel Markov chains, initialized by likelihood weighting, that resample each non-evidence variable from its distribution given its Markov blanket,
  $P(x_i \mid \mathrm{mb}(X_i)) \propto P(x_i \mid \Pa(X_i)) \prod_{C \in \mathrm{Ch}(X_i)} P(c \mid \Pa(C))$,
  and discards a burn-in of 100 sweeps.
\end{itemize}

\begin{algorithm}[H]
  \caption{Likelihood weighting for $P(Q \mid \bm{e})$}
  \label{alg:lw}
  \KwIn{network, query $Q$, evidence $\bm{e}$, number of samples $N$}
  $W[q] \leftarrow 0$ for every state $q$ of $Q$\;
  \For{$s = 1, \dots, N$}{
    $w \leftarrow 1$\;
    \For{$X_i$ in topological order}{
      \eIf{$X_i \in \mathrm{vars}(\bm{e})$}{
        $x_i \leftarrow e_i$;\quad $w \leftarrow w \cdot P(e_i \mid \pa_i)$\;
      }{
        sample $x_i \sim P(X_i \mid \pa_i)$\;
      }
    }
    $W[x_Q] \leftarrow W[x_Q] + w$\;
  }
  \Return $W / \sum_q W[q]$\;
\end{algorithm}

\subsubsection{Classification with a Bayesian network}
For every test applicant, the network computes $P(\text{default} \mid \text{observed features})$ exactly by variable elimination.
Rows with identical evidence share one inference call.
The decision threshold is tuned for F2 on a validation split, exactly as for the other model families, so that all three families are compared under the same rule.
